\newcommand{\titulus}{\nomenFesti{S. Maximiliani Mariæ Kolbe, Presbyteri et Martyris.}
\dies{Die 14. Augusti.}}
\newcommand{\oratio}{\pars{Oratio.}

\noindent Deus, qui sanctum Maximiliánum Maríam, presbýterum et mártyrem, amóre Vírginis Immaculátæ succénsum, animárum zelo et próximi dilectióne replevísti, concéde propítius, ut, eo intercedénte, pro tua glória in servítio hóminum strénue laborántes, usque ad mortem Fílio tuo conformári valeámus.

\pars{Pro pace in universo mundo.} \scriptura{Sir. 50, 25; 2 Esdr. 4, 20; \textbf{H416}}

\vspace{-4mm}

\antiphona{II D}{temporalia/ant-dapacemdomine.gtex}

\vfill

\noindent Deus, a quo sancta desidéria, recta consília et iusta sunt ópera: da servis tuis illam, quam mundus dare non potest, pacem; ut et corda nostra mandátis tuis dédita, et hóstium subláta formídine, témpora sint tua protectióne tranquílla.

\noindent Per Dóminum nostrum Iesum Christum, Fílium tuum, qui tecum vivit et regnat in unitáte Spíritus Sancti, Deus, per ómnia sǽcula sæculórum.

\noindent \Rbardot{} Amen.}
\newcommand{\invitatorium}{\pars{Invitatorium.}

\vspace{-4mm}

\antiphona{E}{temporalia/inv-regemmartyrumsimplex.gtex}}
\newcommand{\hymnusmatutinum}{\pars{Hymnus}

\cuminitiali{I}{temporalia/hym-BeateMartyr.gtex}}
\newcommand{\matversus}{\noindent \Vbardot{} Dómine, véritas tua usque ad nubes.

\noindent \Rbardot{} Iudícia tua abýssus multa.}
\newcommand{\lectioi}{\pars{Lectio I.} \scriptura{Eccl. 8, 5-17}

\noindent De libro Ecclesiástes.

\noindent Qui custódit præcéptum, non experiétur quidquam mali; tempus et iudícium cor sapiéntis intéllegit. Omni enim negótio tempus est et iudícium, et multa hóminis afflíctio; ignórat enim quid futúrum sit, nam quómodo sit futúrum, quis nuntiábit ei? Non est in hóminis potestáte dominári super spíritum nec cohibére spíritum, nec habet potestátem supra diem mortis, nec ulla remíssio est ingruénte bello, neque salvábit impíetas ímpium. Omnia hæc considerávi et dedi cor meum cunctis opéribus, quæ fiunt sub sole, quo témpore dominátur homo hómini in malum suum.

\noindent Et ita vidi ímpios sepúltos, discedéntes de loco sancto; in obliviónem cádere in civitáte, quod ita egérunt: sed et hoc vánitas est. Etenim, quia non profértur cito senténtia contra ópera mala, ídeo cor filiórum hóminum replétur, ut pérpetrent mala. Nam peccátor cénties facit malum et prolóngat sibi dies; verúmtamen novi quod erit bonum timéntibus Deum, qui veréntur fáciem eius. Non sit bonum ímpio, nec prolongábit dies suos quasi umbram, qui non timet fáciem Dómini.

\noindent Est vánitas, quæ fit super terram: sunt iusti, quibus mala provéniunt, quasi ópera égerint impiórum, et sunt ímpii, quibus bona provéniunt, quasi iustórum facta hábeant; sed et hoc vaníssimum iúdico.

\noindent Laudávi ígitur lætítiam quod non esset hómini bonum sub sole, nisi quod coméderet et bíberet atque gaudéret et hoc solum secum auférret de labóre suo in diébus vitæ suæ, quos dedit ei Deus sub sole.

\noindent Cum appósui cor meum, ut scirem sapiéntiam et intellégerem occupatiónem, quæ versátur in terra, quod diébus et nóctibus somnum non capit óculis, ecce intelléxi quod ómnium óperum Dei nullam possit homo inveníre ratiónem eórum, quæ fiunt sub sole; et quanto plus laboráverit homo ad quæréndum, tanto minus invéniet; etiámsi díxerit sápiens se nosse, non póterit reperíre.}
\newcommand{\responsoriumi}{\pars{Responsorium 1.} \scriptura{\Rbardot{} 3 Reg. 8, 39 \& Ps. 138, 16 \Vbardot{} 1 Paral. 28, 9; \textbf{H400}}

\vspace{-5mm}

\responsorium{II}{temporalia/resp-quaesuntincorde-CROCHU-sinedox.gtex}{}}
\newcommand{\lectioii}{\pars{Lectio II.} \scriptura{Eccl. 9, 1-10}

\noindent Omnia hæc cóntuli in corde meo, ut curióse intellégerem quod iusti atque sapiéntes et ópera eórum sunt in manu Dei. Utrum amor sit an ódium, omníno nescit homo: coram illis ómnia. Sicut ómnibus sors una: iusto et ímpio, bono et malo, mundo et immúndo, immolánti víctimas et non immolánti. Sicut bonus sic et peccátor; ut qui iurat, ita et ille qui iuraméntum timet.

\noindent Hoc est péssimum inter ómnia, quæ sub sole fiunt, quia sors éadem cunctis; unde et corda filiórum hóminum impléntur malítia et stultítia in vita sua, et novíssima eórum apud mórtuos. Qui enim sociátur ómnibus vivéntibus, habet fidúciam: mélior est canis vivus leóne mórtuo. Vivéntes enim sciunt se esse moritúros; mórtui vero nihil novérunt ámplius nec habent ultra mercédem, quia oblivióni trádita est memória eórum. Amor quoque eórum et ódium et invídiæ simul periérunt, nec iam habent partem in hoc sǽculo et in ópere, quod sub sole géritur.

\noindent Vade ergo et cómede in lætítia panem tuum et bibe cum gáudio vinum tuum, étenim iam diu placuérunt Deo ópera tua. Omni témpore sint vestiménta tua cándida, et óleum de cápite tuo non defíciat. Perfrúere vita cum uxóre, quam díligis, cunctis diébus vitæ instabilitátis tuæ, qui dati sunt tibi sub sole omni témpore vanitátis tuæ: hæc est enim pars in vita et in labóre tuo, quo labóras sub sole. Quodcúmque fácere potest manus tua, instánter operáre, quia nec opus nec rátio nec sapiéntia nec sciéntia erunt apud ínferos, quo tu próperas.}
\newcommand{\responsoriumii}{\pars{Responsorium 2.} \scriptura{\Rbardot{} Sap. 17, 1 \Vbardot{} ibid. 10, 18; \textbf{H400}}

\vspace{-5mm}

\responsorium{III}{temporalia/resp-magnaenimsunt-CROCHU.gtex}{}}
\newcommand{\lectioiii}{\pars{Lectio III.} \scriptura{O. Joachim Roman Bar, O.F.M. Conv., ed., Wybór Pism, Warszawa 1973, 41-42; 226}

\noindent Ex Epístolis sancti Maximiliáni Maríæ Kolbe.

\noindent {\color{gray} Magno affícior gáudio, dilécte frater, ob flagrántem zelum qui ad promovéndam Dei glóriam te adúrget. Nostris enim tempóribus, non sine mœróre, vidémus epidémicum quemdam morbum, quem \emph{indifferentísmum} vocant, nedum inter sæculáres sed et inter religiósos sodáles, váriis quidem formis propagári. At vero, cum Deus sit dignus infiníta glória, id primum ac potíssimum nobis máximam est, pro tenuitátis nostræ facultáte, eídem glóriam reférre, etsi tantam quam a nobis ipse merétur, éxiles cum simus creatúræ, numquam tribúere valeámus.

\noindent Cum autem Dei glória in animárum salúte, quas Christus próprio sánguine redémit, quam máxime respléndeat, áltius ac præcípuum apostólicæ vitæ stúdium illud esto, quamplúrium animárum salútem immo et altiórem sanctificatiónem procuráre. Quæ vero in hunc finem áptior præsto sit via, scílicet ad divínam comparándam glóriam necnon plurimárum animárum sanctificatiónem, paucis dicam. Deus, qui infiníta sciéntia ac sapiéntia est, adeóque quid a nobis contínuo præstándum sit ut eius adaugeátur glória óptime novit, voluntátem suam nobis per suos in terris vices geréntes, ut plúrimum, maniféstat.

\noindent Est ítaque obœdiéntia, eáque sola, quæ nobis divínam voluntátem certo maniféstat. Póterit quidem supérior in errórem incúrrere, fíeri autem non póterit ut nos, obœdiéntiam sectántes, in errórem inducámur. Tunc tantum obœdiéndi darétur excéptio, si supérior quídpiam præcíperet quod divínæ legis violatiónem, vel in mínimo, manifésto invólveret: quo in casu fidélis ipse Dei voluntátis intérpres non esset.

\noindent Deus isque solus infinítus, sapientíssimus, sanctíssimus, clementíssimus est Dóminus, Creátor et Pater noster, princípium et finis, sapiéntia et poténtia et amor, ómnia ipse Deus. Quidquid vero extra Deum invenítur, eátenus valet quátenus ad ipsum refértur, qui est rerum ómnium Cónditor hominúmque Redémptor, últimus totíus creatiónis finis. Ipse ítaque est, qui adorábilem suam voluntátem per suos in terris vices geréntes nobis maniféstat et nos ad se trahit, álias quoque ánimas per nos ad se tráhere ac perfectióre sibi caritáte coniúngere inténdens.}

\noindent Vide, frater, quanta sit, per Dei misericórdiam, condiciónis nostræ dígnitas. Per obœdiéntiam, exiguitátis nostræ términos véluti transcéndimus, divinǽque voluntáti conformámur, quæ infiníta sua sapiéntia ac prudéntia ad recte agéndum nos dírigit. Immo, eídem divínæ inhæréntes voluntáti, cui nulla res creáta resístere potest, cunctis fortióres effícimur.

\noindent Hæc est sapiéntiæ ac prudéntiæ sémita, hæc única via est, qua Deo summam glóriam reférre valeámus. Si ália et áptior via esset, eam Christus verbo exemplóque suo nobis profécto manifestásset. Sed diutúrnam eius in Názareth vitæ consuetúdinem divína Scriptúra iis perstrínxit verbis: Et erat súbditus illis, reliquúmque vitæ cursum sub obœdiéntiæ véluti signo nobis adumbrávit, passim osténdens eum in terras descendísse ut fáceret voluntátem Patris.

\noindent Diligámus ítaque, fratres, summópere diligámus cæléstem amantíssimum Patrem, sitque huius perféctæ caritátis arguméntum obœdiéntia nostra, tunc vel máxime exercénda cum própriæ voluntátis sacrifícium nobis postulétur. Nec enim sublimiórem librum nóvimus, quam Iesum Christum crucifíxum, ut in amóre Dei progrediámur.

\noindent Quæ ómnia expedítius per Vírginem Immaculátam obtinébimus, cui Deus benigníssimus suæ dispensatiónem misericórdiæ commísit. Nullum équidem dúbium, quin Maríæ volúntas ipsíssima Dei volúntas nobis sit. Cumque nosmetípsos eídem dedicámus, instruménta in eius mánibus, quemádmodum et ipsa in Dei mánibus, divínæ misericórdiæ effícimur. Sinámus ítaque nos ab ea dírigi, sinámus ab ea manudúci, simúsque sub eius ductu quiéti ac secúri: cuncta enim nobis ipsa prospíciet, cuncta providébit, córporis animíque necessitátibus prompta succúrret, difficultátes et angústias ipsa removébit.}
\newcommand{\responsoriumiii}{\pars{Responsorium 3.} \scriptura{\textbf{H373}}

\vspace{-5mm}

\responsorium{III}{temporalia/resp-istecognovit-CROCHU-cumdox.gtex}{}}
\newcommand{\hymnuslaudes}{\pars{Hymnus}

\cuminitiali{VI}{temporalia/hym-MartyrDei.gtex}}
\newcommand{\lectiobrevis}{\pars{Lectio Brevis.} \scriptura{2 Cor. 1, 3-5}

\noindent Benedíctus Deus et Pater Dómini nostri Iesu Christi, Pater misericordiárum et Deus totíus consolatiónis, qui consolátur nos in omni tribulatióne nostra, ut possímus et ipsi consolári eos, qui in omni pressúra sunt, per exhortatiónem, qua exhortámur et ipsi a Deo; quóniam, sicut abúndant passiónes Christi in nobis, ita per Christum abúndat et consolátio nostra.}
\newcommand{\responsoriumbreve}{\pars{Responsorium breve.} \scriptura{Ex. 15, 2}

\cuminitiali{VI}{temporalia/resp-fortitudomeaetlausmea.gtex}}
\newcommand{\benedictus}{\pars{Canticum Zachariæ.} \scriptura{Io. 15, 13; \textbf{H363}}

\vspace{-4mm}

\antiphona{I g}{temporalia/ant-majoremcaritatem.gtex}

\vspace{-2mm}

\scriptura{Lc. 1, 68-79}

\vspace{-2mm}

\cantusSineNeumas
\initiumpsalmi{temporalia/benedictus-initium-i-g-auto.gtex}

%\vspace{-1.5mm}

\input{temporalia/benedictus-i-g.tex} \Abardot{}}
\newcommand{\preces}{\noindent Fratres, Salvatórem nostrum,~\gredagger{} testem fidélem, per mártyres interféctos propter verbum Dei, celebrémus,~\grestar{} clamántes:

\Rbardot{} Redemísti nos Deo in sánguine tuo.

\noindent Per mártyres tuos,~\gredagger{} qui líbere mortem in testimónium fídei sunt ampléxi,~\grestar{} da nobis, Dómine, veram spíritus libertátem.

\Rbardot{} Redemísti nos Deo in sánguine tuo.

\noindent Per mártyres tuos,~\gredagger{} qui fidem usque ad sánguinem sunt conféssi,~\grestar{} da nobis, Dómine, puritátem fideíque constántiam.

\Rbardot{} Redemísti nos Deo in sánguine tuo.

\noindent Per mártyres tuos,~\gredagger{} qui, sustinéntes crucem, tua vestígia sunt secúti,~\grestar{} da nobis, Dómine, ærúmnas vitæ fórtiter sustinére.

\Rbardot{} Redemísti nos Deo in sánguine tuo.

\noindent Per mártyres tuos,~\gredagger{} qui stolas suas lavérunt in sánguine Agni,~\grestar{} da nobis, Dómine, omnes insídias carnis mundíque devíncere.

\Rbardot{} Redemísti nos Deo in sánguine tuo.}
\newcommand{\benedicamuslaudes}{\cuminitiali{}{temporalia/benedicamus-memoria-laudes.gtex}}
\include{hebdomadaxix}
\include{feriavi}
